\documentclass[ngerman,a4paper,12pt,headsepline]{scrreprt}

%----------------- PDF CONFIG ----------------- %
\pdfinfo{    
     /Title (PDF-Titel) 
     /Subject   (PDF-Thema)    
     /Author  (Vorname Nachname) 
     /Keywords   (Stichwort1,Stichwort2)      
} 

\title{Testframework}
\author{Can Kedik}
\date{\today}



%----------------- PAKETE INKLUDIEREN ----------------- %

\usepackage{geometry} % Packet für Seitenrandabständex und Einstellung für Seitenränder
\usepackage[ngerman]{babel} % deutsche Silbentrennung

\usepackage{booktabs} %entzerrt die Tabellenzeilen und bietet verschieden dicke Unterteilungslinien
\usepackage{longtable}% Tabellen können sich nicht über mehrere Seiten 
\usepackage{graphicx} % kann LaTeX Grafiken einbinden

\usepackage{setspace}
\usepackage{mathabx} 
\usepackage[applemac]{inputenc} % Umlaute unter Mac werden automatisch gesetzt
\usepackage[T1]{fontenc} % Zeichenencoding
\usepackage{lmodern} % typographische Qualität 
\frenchspacing % Schaltet den zusätzlichen Zwischenraum ab
\usepackage{fix-cm}
\usepackage{varioref}
\usepackage{hyperref} % verwandelt alle Kapitelüberschriften, Verweise aufs Literaturverzeichnis und andere Querverweise in PDF-Hyperlinks
\usepackage{cleveref}
\usepackage{color}
\usepackage{url}
\usepackage[babel, german=quotes]{csquotes}
\usepackage{array}
\usepackage{siunitx}
\usepackage[backend=biber,style=alphabetic]{biblatex}
\bibliography{literatur}


\usepackage[nottoc]{tocbibind}



% für Listings
\usepackage{listings}
\lstset{numbers=left, numberstyle=\tiny, numbersep=5pt, stepnumber=4, keywordstyle=\color{black}\bfseries\itshape, stringstyle=\ttfamily,showstringspaces=false,basicstyle=\footnotesize,captionpos=b}
\lstset{language=java}



%----------------- FARBEN DEFINIEREN ----------------- %
\definecolor{gray}{gray}{0.95} % Listingsbackground

%----------------- LAYOUT SETZEN ----------------- %
\geometry{left=2.5cm, right=2cm, top=2.5cm, bottom=2cm}
\linespread {1.25}\selectfont %1.25 da er von Haus aus 1.2 ist und 1,25 * 1,2 = 1,5 isch



%-------##-------##-------##------- ANFANG INHALT -------##-------##-------##-------%
\begin{document}



\pagenumbering{roman} % Seitennummer

%----------------- DECKBLATT -----------------%
 %----------------- KONFIGURATION ----------------- %
\pagestyle{empty} % enthalten keinerlei Kopf oder Fuß 

%----------------- HDA FBI Logo ----------------- %
\begin{figure}[t]
	\centering
	\includegraphics[width=0.6\textwidth]{Abb/logo_fbi}
\end{figure}

%----------------- INHALT ----------------- %

\begin{center}
\Large Hochschule Darmstadt \\
\normalsize \textsc{- Fachbereich Informatik -} \\

% Whitespace
\vspace{105 pt}

\Huge Name der Arbeit \\ 
\normalsize
\vspace{20 pt}

Abschlussarbeit zur Erlangung des akademischen Grades \\ 
Bachelor of Science (B.Sc.) 

\vspace{75 pt}


vorgelegt von \\
\vspace{5 pt}
Can Kedik \\
731620
\vspace{115 pt}

\begin{tabular}[h]{p{4cm}l l}
	Referent: & Prof. Dr. Ronald C. Moore\\
	Korreferent: & Prof. Dr. Bettina Harriehausen-M"uhlbauer
\end{tabular}


\end{center}

 
%----------------- ABSTRACT -----------------%
 %----------------- KONFIGURATION ----------------- %
\pagestyle{empty} % enthalten keinerlei Kopf oder Fuß


\chapter*{Abstract} % (fold)
\label{cha:abtract}
	Lorem ipsum dolor sit amet, consetetur sadipscing elitr, sed diam nonumy eirmod tempor invidunt ut labore et dolore magna aliquyam erat, sed diam voluptua. At vero eos et accusam et justo duo dolores et ea rebum. Stet clita kasd gubergren, no sea takimata sanctus est Lorem ipsum dolor sit amet. Lorem ipsum dolor sit amet, consetetur sadipscing elitr, sed diam nonumy eirmod tempor invidunt ut labore et dolore magna aliquyam erat, sed diam voluptua. At vero eos et accusam et justo duo dolores et ea rebum. Stet clita kasd gubergren, no sea takimata sanctus est Lorem ipsum dolor sit amet.

\begin{quote}
Dies ist ein Zitat.
\end{quote}

verstand, scheinen nun doch vorueber zu Dies ist der Text sein. \\
siehe: http://janeden.net/die-praeambel

 
 
%----------------- VERZEICHNISSE -----------------%
\tableofcontents % Inhaltverzeichnis

% ----- Abbildungen ----- %
%\addcontentsline{toc}{section}{Abbildungsverzeichnis} % falls in Inhalsverzeichnis
\listoffigures

% ----- Tabellen----- %
% \addcontentsline{toc}{section}{Tabellenverzeichnis}  % falls in Inhalsverzeichnis
% \fancyhead[L]{Abbildungsverzeichnis / Abkürzungsverzeichnis} %Kopfzeile links
\listoftables

% ----- Listings ----- %
% Listingverzeichnis soll im Inhaltsverzeichnis auftauchen
% \addcontentsline{toc}{section}{Listingverzeichnis}
% \fancyhead[L]{Abbildungs- / Tabellen- / Listingverzeichnis} %Kopfzeile links
\renewcommand{\lstlistlistingname}{Listingverzeichnis}
\lstlistoflistings

\pagestyle{headings} % zurueck setzen von roemische seitenanzahl


%----------------- KAPITEL : EINLEITUNG  ----------------- %
\pagenumbering{arabic}
\chapter{Einf"uhrung}
\label{cha:einfuehrung}

Lorem ipsum dolor sit amet, consetetur sadipscing elitr, sed diam nonumy eirmod tempor invidunt ut labore et dolore magna aliquyam erat, sed diam voluptua. At vero eos et accusam et justo duo dolores et ea rebum. Stet clita kasd gubergren, no sea takimata sanctus est Lorem ipsum dolor sit amet. Lorem ipsum dolor sit amet, consetetur sadipscing elitr, sed diam nonumy eirmod tempor invidunt ut labore et dolore magna aliquyam erat, sed diam voluptua. At vero eos et accusam et justo duo dolores et ea rebum. Stet clita kasd gubergren, no sea takimata sanctus est Lorem ipsum dolor sit amet.




%----------------- KAPITEL : FRAMEWORK  ----------------- %
\chapter{Das Framework}
\label{cha:framework}
\section{Verwendete Hardware} \label{sec:verwendeteHardware}
Die in \cref{sec:testbeschreibung} beschriebenen Versuche wurden mit folgender Hardware durchgef"uhrt. %
\begin{itemize}
\item Zwei Switches
\item Vier Raspberry Pis der ersten Generation
\item Ein Raspberry Pi der zweiten Generation
\item Mehrere Ethernet Kabel
\end{itemize}
Diese Ger"ate wurden in einem eigenst"andigend Netzwerk zusammengeschaltet. So sind an jedem Switch zwei Pis %
der ersten Generation angeschlossen w"ahrend an einem der Switches der Pi der zweiten Generation angeschlossen %
ist (siehe \cref{fig:AufbauVomNetzwerk}). Welche Software auf den Pis verwendet wurde, wird in \cref{sec:aufbauSoftware} erkl"art. 

\section{Aufbau der Software} \label{sec:aufbauSoftware}
Das Testframework besteht aus drei verschiedenen Teilen. %
\begin{itemize}
\item Zabbix %
\item Eigentwicklung auf dem Server %
\item Eigentwicklung auf dem Agent %
\end{itemize}
Die Eigentwicklungen sind alle in Bash programmiert, w"ahrend Zabbix eine bereits fertige Open Source L"osung %
ist. In den folgenden Abschnitten werde ich einen Einblick in diese Teile geben. %

\subsection{Zabbix}
Zabbix ist ein Open Source Netzwerk Monitoring System. Die erste Version wurde von Alexei Vladishev entwickelt, welches %
seit 2005 von der Firma Zabbix SIA weiterentwickelt wird \autocite{zabbix:Web}. % 
Ein weiterer bekannter Vertreter der Netzwerk Monitor Systeme ist Nagios, welches wie Zabbix unter der GPL Lizenz vertrieben wird \autocite{wiki:Nagios}. % 
Beide Systeme basieren auf einer Client-Server Architektur. Im weiteren wird jedoch nur Zabbix betrachtet. %
Zabbix besteht aus zwei Komponenten. %

\begin{description}
\item[Zabbix Server]Der Server hat eine auf PHP basierende Weboberfl"ache "uber die es f"ur den Benutzer m"oglich ist die Agents zu %
konfigurieren. So k"onnen manuell die Templates erstellt werden, die den Zabbix Agents mitteilen, welche Informationen dem %
Server zu "ubermitteln sind. %
"Uber die API Schnittstelle ist es m"oglich Programme f"ur den Server zu schreiben und diese auszuf"uhren. %
Der Zabbix Server ist ein sich selbst "uberwachender Agent. Dieser ist jedoch nicht beteiligt %
im allgemeinen Prozess des Frameworks. %
Trotzdem wird in den folgenden Kapiteln der Zabbix Server nicht aussen vor gelassen und in die %
den Tests betrachtet werden, dabei werden haupts"achlich die von Zabbix generierten Werte zur Perfomance und zum %
Traffic auf dem Ethernet Port genommen. Der Wert Perfomance betrachtet zwei Dinge, einmal die sogenannte Queue, welche %
angibt wie viele von den Agents "ubermittelten Werten darauf warten vom Zabbix Server verarbeitet zu werden und die %
die verarbeiteten Werte pro Sekunde. Diese Werte lassen einen Schluss auf die Leistung des Servers schlie"sen. %
Sollte die Queue einen bestimmten schwellwert erreichen wird auf dem Zabbix Dashboard eine Warnung ausgegeben das der %
Server mit der Verarbeitung nicht hinterherkommt. Dies kann soweit gehen das die gesammelten Daten auf dem Server nicht %
aktuell oder sogar fehlerhaft sind. Jedoch ist dieses Szenario in diesem Framework unwahrscheinlich, da die gr"o"se des %
Frameworks "uberschaubbar ist. Der Traffic auf dem Ethernet Port Eth0 wird betrachtet, da alle von den Agents gesammelten Informationen %
"uber diesen an den Server gelangen. Dabei wird der eingehende sowie auch der ausgehende Traffic betrachtet. 

\item[Zabbix Agent]Die Clients, die im Netzwerk "uberwacht werden sollen sind die sogenannten Agents, die an den Server %
die Informationen weiterleiten, die vom Server gefordert werden. In den sp"ateren Kapiteln wird der Hauptfokus auf der %
I/O--Last der Festplatte und wie beim Server dem ein/- und ausgehenden Traffic auf dem Ethernet Port Eth0 liegen. %
\end{description} 

\subsection{Eigenentwicklung}
Die selbstentwickelte Software wird in zwei Kategorien unterteilt; ein Teil der Software l"auft %
auf den Agents, diese haben den Zweck Netzwerk Traffic zu erzeugen. Dadurch entsteht auf dem Netzwerk und %
auf den Zabbix Agents eine Nutzlast, die vom Zabbix Server gesammelt werden kann. Der zweite Teil der Software %
die auf dem Server l"auft, die Software auf dem Server unterst"utzt den Entwicklungsprozess auf den Agents. %
\begin{enumerate}
\item Zabbix Server 
\begin{description}
\item[Update Script:]Dieses Script wird von der Software Entwicklungsmaschine ausgef"uhrt. Es aktualisiert den Code auf den Endger"aten, %
die die Programme Hintergrundrauschen, Pinger, Synchronize und Startrauschen aktualisieren. Dabei wird mittels Secure Copy %
der Quellcode auf das Endger"at gespielt.

\item[Get logs:]Die Skripte Pinger und Hintergrundrauschen erstellen jeweils auf den Endger"aten logfiles. %
Da es jedoch ein sehr hoher Verwaltungsaufwand w"are auf den Endger"aten die Logfiles weiterzuverwerten, %
werden mit dem Skript Get Logs die auf den Agents gelagerten Logfiles auf dem Rechner der dieses Skript %
startet gesammelt. Somit hat man die von den Endger"aten gesammelten Daten auf einem Rechner und kann %
dessen Weiterverarbeitung betreiben.      

\item[Calculate Average:]Berechnet den Durschnitt der Dauer, die ein Paket braucht, um erfolgreich versendet %
zu werden.
 
\end{description}

\item Zabbix Agent

\begin{description}
\item[Hintergrundrauschen:]Diese Eigentwicklung stellt mit Zabbix die Kernkomponente des Frameworks dar. Wie %
der Rest der selbstentwickelten Software wurde sie komplett in Bash programmiert, da das Framework ausschlie"slich %
in einer Linux Umgebung entwickelt, getestet und verwendet wird. Hintergrundrauschen schickt "uber Secure Copy, %
Pakete von einem Agent zum anderen. So wird eine %
Last auf dem Netzwerk erzeugt, die Mithilfe des Zabbix Servers gemessen werden kann. Ausserdem speichert Hintergrundrauschen %
die Dauer, bis ein Paket erfolgreich bei seinem zuf"allig ausgew"ahlten Empf"anger angekommen ist. %

\item[Startrauschen:]wird automatisch zum Start eines der Endger"ate ausgef"uhrt. Es dient als eine Zeitschaltuhr %
und erm"oglicht ein zeitversetztes Starten des Scripts Hintergrundrauschen, wodurch man den sequentiellen %
Anstieg an Last im Netzwerk beobachten kann.

\item[Synchronize:]Raspberry Pis besitzen keine eigene Batterie wie es handels"ubliche Rechner haben, %
deshalb ist nach jedem Neustart die Uhrzeit der Pis unzuverl"assig. Dieses Programm synchronisiert %
die Uhrzeiten der Agents mit der Zeit des Servers. %

\item[Pinger:]wird zusammen mit dem Hintergrundrauschen ausgef"uhrt und ist um den Linux eigenen Ping Befehl %
aufgebaut. Mittels Pinger wird die Latenz unter den einzelnen Endger"aten gemessen. %
\end{description}
\end{enumerate}
\section{Einsatz im Netzwerk} 

Mit dem Einsatz der im vorherigen Abschnitt vorgestellten Software ist das Testframework aufgebaut. In der \cref{fig:AufbauVomNetzwerk} %
wird dargestellt, wie die einzelnen Komponenten zusammenspielen. Hier sieht man, dass an einem Switch zwei Raspberry Pis %
und an einem anderen Switch drei Raspberry Pis angeschlossen sind. Einer von diesen Pis ist der Zabbix Server, der %
die Aktiven Hosts im Netzwerk "uberwacht. 

\label{sec:einsatzImNetzwerk}
\begin{figure}[htbp]
\centering
\includegraphics*[width=0.9\linewidth]{Abb/Netzschaltung3}

\caption{Aufbau des Netzwerks}
\label{fig:AufbauVomNetzwerk}
\end{figure}















%----------------- KAPITEL : TESTAUFBAU  ----------------- %
\chapter{Zabbix}
\label{cha:zabbix}
In \cref{cha:framework} wurde schon die verwendete Netzwerk Monitoring Software angeschnitten. In diesem Abschnitt wird nun %
ein tieferer Einblick in den Aufbau und die Verwendung von Zabbix gew"ahrt. Zabbix erm"oglicht es %
auch das Monitoring als ein Verteiltes Monitoring aufzuziehen mit mehreren Servern, diese sogenannten Proxies %
sind jedoch kein Bestandteil dieser Thesis und werden deshalb nicht weiter betrachtet.

\section{Zabbix Server}
\label{sec:server}
Der Zabbix Server ist die Zentrale Komponente des vorgestellten Frameworks. Dem Server werden von den Agents und Proxies %
Informationen zugeschickt. Die Aufgabe des Server ist es die Daten zu speichern und zu verarbeiten. %
Der Server speichert auch die Konfiguration der einzelnen Agents die sich im Netzwerk befinden. %
Um die Konfiguration der Agents einstellen zu k"onnen m"ussen jedoch erst einmal die Agents im Zabbix %
Server registriert werden. Die dazu ben"otigten Daten sind die Agent IP und ein eindeutiger Name. %
Wenn nun der Agent im Server registriert worden ist kann man im Zabbix Server den Agents verschiedene Templates auflegen. %
Das in dem in \cref{cha:framework} gezeigte Framework verwendet prim"ar das Template: \emph{Template OS %
Linux}. Dieses Templates ist eines von vielen vordefinierten Templates. Es ist auch m"oglich selber Templates %
zu erstellen, Templates sind in einem XML Format abgespeichert und k"onnen dadurch einfach unter Nutzern von Zabbix % 
ausgetauscht werden. Die Firma die den vertrieb von Zabbix regelt hat extra daf"ur die Zabbix Share eingef"uhrt %
auf dem Zabbix Templates und vieles mehr ausgetauscht werden kann \autocite{zabbix:share}. \ \\ %
Templates enthalten sogenannte Items die "uber Active Checks oder Passive Checks Informationen von %
einem vorher im Zabbix Server registrierten Agent anfordern. Ein Passive Check geschieht "uber einen %
Request diese sind in JSON verfasst und werden in geb"undelt verschickt um Bandbreite zu sparen. %
Der Prozess beeinhaltet f"unf Schritte: %
\begin{enumerate}
\item Der Server "offnet eine TCP Verbindung %
\item Der Server sendet einen request als Beispiel agent.version %
\item Der Agent empf"angt den request und antwortet mit der Versionsnummer %
\item Der Server verarbeitet die Antwort und erh"alt als Ergebniss Zabbix3.0.0rc %
\item Die TCP Verbindung wird geschlossen %
\end{enumerate}
\autocite{zabbix:activepassive}. Passive Checks und Active Checks unterscheiden sich darin, wer den request %
ausl"ost. Bei einem passiv Check sind die Hosts selber inaktiv bis vom Server ein request eintrifft, bei Active %
Checks wird der Agent selber Aktiv. Der Agent selber sendet dann an den Server einen Request in dem der Agent nach den %
Daten fragt die der Server erwartet sind\autocite{zabbix:activepassive2}. %
Dies ist der Fall wenn vom Agent Daten abgefragt werden sollen die nicht vom %
Betriebssystem des Agents gesammelt werden. Das bedeutet das der Agent selber nun die Daten die der Server erwartet %
sammeln muss, dies geschieht meistens "uber externe Programme, sollte ein Agent zu viele Active Checks haben kann es %
sich auf die Performanz des Hosts auswirken. \ \\ % 

"Uber die API ist es m"oglich Programme f"ur den Server zu schreiben mit der man zum Beispiel %
eine eigene Benutzeroberfl"ache erstellen kann um die Konfigurationen auf dem Server zu verwalten. %
"Uber die Api ist es auch m"oglich einen eigenen Zugriff auf die dem Server zugrunde liegende %
Datenbank herzustellen. %
Der Zabbix Server basiert auf einem Apache Web Server. Um das Zabbix Frontend nutzen zu k"onnen wird %
PHP 5.4.0 ben"otigt, jedoch funktioniert PHP 7 noch nicht. Die Daten und Statistiken die Zabbix sammelt %
werden in einer Datenbank gespeichert. Es werden f"unf verschiedene Datenbanken unterst"utzt. %
\begin{itemize}
\item MySql Version 5.0.3 aufw"arts
\item Oracle Version 10g aufw"arts
\item PostgreSQL Version 8.1 aufw"arts
\item SQLite Version 3.3.5 aufw"arts
\item IBM DB2 Version 9.7 aufw"arts (Noch nicht fehlerfrei)
\end{itemize}
\autocite{zabbix:req}
Der Zabbix Server selber ist auch als ein Agent konfiguriert, der sich selber "uberwacht. Jedoch ist der Server %
nicht im allgemeinen Prozess des in \cref{cha:framework} vorgestellten Frameworks t"atig. %
Trotzdem wird in den folgenden Kapiteln der Zabbix Server nicht aussen vor gelassen und in die %
den Tests betrachtet werden, dabei werden haupts"achlich die von Zabbix Statistiken zur Perfomance und zum %
Traffic auf dem Ethernet Port des Servers genommen. Der Wert Perfomance betrachtet zwei Dinge, einmal die sogenannte Queue, welche %
angibt wie viele von den Agents "ubermittelten Werten darauf warten vom Zabbix Server verarbeitet zu werden und die %
die verarbeiteten Werte pro Sekunde. Diese Werte lassen einen Schluss auf die Leistungsanforderungen des Servers schlie"sen. %
Sollte die Queue einen bestimmten schwellwert erreichen wird auf dem Zabbix Frontend eine Warnung ausgegeben das der %
Server mit der Verarbeitung nicht hinterherkommt. Dies kann soweit gehen das die gesammelten Daten auf dem Server fehlerhaft %
sind. Jedoch ist dieses Szenario in diesem Framework unwahrscheinlich, da die gr"o"se des %
Frameworks "uberschaubar ist. Der Traffic auf dem Ethernet Port Eth0 wird betrachtet, da alle von den Agents gesammelten Informationen %
"uber diesen an den Server gelangen. Dabei wird der eingehende sowie auch der ausgehende Traffic betrachtet.

\section{Zabbix Agent}
\label{sec:agent}
Der Zabbix Agent wird auf dem zu "uberwachenden System installiert. Der Agent sammelt die Daten "uber im Betriebssystem %
integrierte Monitoring Funktionen oder "uber die Zabbix eigene API und sendet diese je nach Art des Checks diese Informationen an den Agent. %
Da der Agent schon im Betriebssystem integrierte Funktionen benutzt ist dieser sehr effizient und verbraucht kaum sp"urbar %
die Ressourcen des Host Systems. Anders als der Server besitzt der Agent kein eigenes Frontend und speichert die Werte nicht %
in einer Datenbank. Der Agent ist so konstruiert das dieser weitesgehend ohne Administrator Rechte funktionieren kann. % 
Da es auch m"oglich ist "uber einen Fernzugriff Agent Hostsysteme auszuschalten oder neuzustarten m"ussen f"ur diese Funktionen %
dem Agent in der Rechteverwaltung diese Rechte zugeteilt werden. Dadurch wird das Sicherheits%
risiko f"ur den Host erheblich reduziert. F"ur gew"ohnlich wird f"ur den Agent ein eigener User angelegt. Die Anwendung des Agents %
l"auft unter Unix Systemen als ein Daemon und unter Windows als ein Systemdienst. Zabbix unterst"utzt jedoch noch mehr %
Betriebssysteme eine Auswahl davon w"aren: 
\begin{itemize}
\item Linux
\item Windows: alle Desktop und Server Versionen seit 2000
\item OpenBSD
\item Mac OS X
\item Solaris 9,10,11
\end{itemize}
\autocite{zabbix:agent}.
Im Aufbau des Frameworks das in dieser Arbeit weiter betrachtet wird laufen die Zabbix Agents unter dem g"angigen Raspberry Pi %
Betriebssystem Raspbian, welches ein Debian Port ist. Um den Host jedoch verwenden zu k"onnen muss man in %
den Konfigurations Dateien die Ports die der Server f"ur das versenden von Requests benutzt eingestellt werden und auch die IP des Servers %
muss eingetragen werden. Da der Agent sonst eingehende Requests verwirft und somit die Sicherheit f"ur den Host gew"ahrleistet. %
In den Konfigurationsdateien der Agents ist es auch m"oglich %
die Active Checks zu notieren, was in diesem Versuchsaufbau auch gemacht wurde. Da das Template OS Linux keine Items hatte die das %
"uberwachen von der Festplatten Last hatte kann man dies "uber die Konfigurationsdateien einstellen. Das sieht in der Konfigurationsdatei %
dann so aus.
\begin{quote}

UserParameter=custom.vfs.dev.read.ops[*],customParaScript.sh \$1 4 

\end{quote}

Wenn dann in der Antwort zu einem Active Check nach dem \emph{custom.vfs.dev.read.ops[*]} gefragt wird, wird das Bash Script customParaScript.sh %
ausgef"uhrt. Die Auszeichnung emph{UserParameter=} gibt an das es sich hierbei um ein selbstdefiniertes Item handelt, welches mithilfe von Zabbix %
"uberwacht werden soll. Das Script customParaScript.sh dient dazu die Last auf der Festplatte nachzupr"ufen %
dazu liest es aus der Datei die in Unix Systemen unter /proc/diskstats liegt, nach den zugeh"origen Werten \$1 und 4 sind Variablen die f"ur die %
verarbeitung notwendig sind. Der in diesem Beispiel genannte UserParameter liest die Lese Operationen pro Sekunde aus der Datei und leitet diese %
an den Server weiter.


%----------------- KAPITEL : VERSUCHE  ----------------- %
\chapter{Versuche}
\label{cha:versuche}
Die beste Art ein Netzwerktestframework zu Testen ist in dem man das Framework benutzt %
dazu habe ich am Anfang der Arbeiten selber mehrere Tests definiert, jeder dieser Tests %
ist eine in einem Netzwerk m"ogliche Fehlerquelle die die Leistung des Netzwerkes %
beeintr"achtigen kann. Um Fehler identifizieren %
zu k"onnen m"ussen jedoch auch die Normalwerte bestimmt werden. Dazu muss erstmal ein %
St"orungsfreies Netzwerk aufgebaut werden und dieses muss beobachtet werden. %
Die in diesem Netz gesammelten Werte stellen die Basis f"ur unsere sp"ateren Ergebnisse dar. %
Der zweite Schritt ist es das Netz mit Fehlern aufzubauen und zu beobachten. Dadurch %
haben wir die Werte eines Fehlerbehafteten Netzwerkes und k"onnen diese nun mit dem %
fehlerfreien Netzwerk vergleichen. Dadurch sollte es uns nun m"oglich sein automatisierte %
Aussagen "uber die Art des Fehlers in einem Netzwerk zu treffen. %

\section{Versuchsbeschreibung}
\label{sec:versuchbeschreibung}
\begin{description}
\item[Ethernetkabel ohne Isolierung] Ethernetkabel werden oft sehr strapaziert %
dies kann dazu fuhren das sich die Isolation l"ost und somit Strahlung einwirkt. %
Dies kann die Leistung im Netzwerk beeintr"achtigen. 
\item[Falsch gedrehtes Twisted Pair Kabel] TP Kabel sind die am meisten in einem Ethnernet %
Netzwerk vorkommenden Kabel, es entsteht eine hohe St"orungsicherheit bei TP Kabeln %
sollten jedoch TP Kabel entdrillt worden sein kann eine hohe Strahlenbelastung auftreten. %
\item[Loop] Ein Loop ist wenn ein Switch mit sich selbst verbunden ist, der Switch sendet %
dann st"andig Datagramme an sich selbst und blockiert so jeden Traffic auf der Leitung. %
\item[Nicht angeschlossenes Kabel] Wie verh"alt sich das Netzwerk wenn ein Kabel von einem %
Endger"at im laufenden Betrieb entkoppelt wird.
\item[Forkbomb] Die Forkbomb ist ein Programm das von sich rekursiv Kopien erstellt %
so das der Computer alle sein Resourcen verbraucht nur noch diesen Code auszufuhren. %
\begin{quote} 
:()\{ :|:\&\};: [https://de.wikipedia.org/wiki/Forkbomb] 
\end{quote}
\item[Festplatte l"auft voll] In einem Netzwerk werden st"andig Daten versand, was also %
passiert wenn die Festplatte von einem Rechner voll"auft und der PC somit Arbeitsunf"ahig wird. %
\item[IP Adresse doppelt belegt] Firmennetze benutzen statische IPs damit Webserver immer unter %
dem selben Namen erreichbar bleiben, das jedoch zu konfigurieren kann zu Fehlern Menschlichen %
Fehlern fuhren, was passiert in einem Netz wenn zwei Rechner sich eine IP teilen. %
\item[Kollisionsdom"anen] Kollisionsdom"anen sind bereiche in einem Netzwerk wo sich mehrere Rechner eine %
eine Leitung teilen, dies ist geschieht in der Physikalischen Ebene eines Netwerks. Die Frage ist kann
man eine Kollisionsdom"ane erkennen. 
\end{description}
\section{Raspberry Pi Versuche}
\label{sec:raspberryPiVersuche}
\subsection{Gemachte Tests}
\subsection{Ergebnisse}

\section{Virtual Machine Versuche}
\label{sec:VMVersuche}
\subsection{Gemachte Tests}
\subsection{Ergebnisse}



%----------------- KAPITEL : VERGLEICH------------ %	
\chapter{Ergebnisse}
\label{cha:ergebnisse}
In diesem Kapitel werden die Ergebnisse aus dem \cref{cha:versuche} miteinander verglichen. %
Dabei wird als erstes die Anzahl der verschickten Pakete und die Menge der verschickten %
Bytes miteinander verglichen und die Zeiten die, ein Paket Paket zum erreichen des %
Zielsystems Ben"otigt. Es wird ein Blick auf den Ping im Netzwerk und die Auslastung auf den %
Ethernet Ports geworfen. Es wird viel Wert auf einen m"oglichst ausgeglichen %
Traffic an den Ports gelegt. Als zweites wird ein Blick auf die Auslastung des Prozessors gelegt. %
Vor allem wird versucht die Paketegr"o"se zu finden, die die CPU am wenigsten in Anspruch nimmt, %
so das der Endnutzer einen m"oglichst Leistungsf"ahigen Host verwenden kann. Als drittes %
und letztes der im \cref{cha:versuche} betrachteten Metriken kommt die Auslastung der Festplatte. %
Bevor es zur Schlussfolgerung kommt wird jedoch noch ein aufgetretener Fehler betrachtet der den %
Host Tusk in \cref{sec:2000MBTest} betroffen hatte. Zum Schluss dieses Kapitels wird eine Schlussfolgerung %
aus den Ergebnissen gezogen. % 

Das Ziel in diesem Kapitel ist es, die Paketgr"o"se zu ermitteln mit der man den h"ochsten Datenaustausch zwischen den %
Hosts erreichen kann, ohne das die Hosts Handlungsunf"ahig werden und Endnutzer noch auf %
den Hosts simple Arbeiten vollbringen k"onnen. Darunter wird gez"ahlt das der Nutzer in der %
Lage ist den Rechner bedienen zu k"onnen und nicht das es auf dem Rechner m"oglich ist aufw"andige Software %
zu verwenden, wie z.B: Photoshop oder "ahnliches.  

\section{Versandte Daten}
\label{sec:verschickteDaten}
Da dieses Framework darauf ausgelegt ist den h"ochsten Datenverkehr zu bestimmen ist %
Auslastung der Ethernet Ports auf den jeweiligen Hosts der erste Indikator um zu bestimmen %
in welcher Paketgr"o"se der h"ochste Austausch an Daten besteht. Dazu werden die durchschnittlichen %
Werte aus der \cref{tab:compTraffic} miteinander verglichen. Diese Durchschnittswerte beziehen sich %
auf die Tabellen aus \cref{cha:versuche} die den Traffic der Ports aufgelistet hatten. Man sieht %
auf den ersten Blick das sich der Traffic sich bei allen drei Tests nur um Kilobyte unterscheidet. %
Bei einer Paketgr"o"se von 200 Megabyte ist jedoch ersichtlich das diese den h"ochsten eingehenden %
Traffic auf den Hosts erzeugt. Bei einer Paketgr"o"se von 2000 Megabyte liegt dieser knapp unter %
8 Megabit/s und ist damit die Paketgr"o"se die den geringsten eingehenden Traffic erzeugt. Bei %
den 20 Megabyte Tests kommen 8,34 Megabit/s durchschnittlich an, dieser liegt im Mittelfeld %
im Vergleich zum eingehenden Traffic. Betrachtet man nun den ausgehenden Traffic ist wieder %
der Test bei dem 200 Megabyte gro"se Pakete verschickt worden sind an der Spitze mit 8,67 %
Megabit/s. Der 2000 Megabyte Test ist diesmal mit 8,17 Megabit/s "uber einen ausgehenden %
Traffic von 8 Megabit/s gekommen bildet aber auch hier wieder das Schlusslicht der "Ubertragung. %
Bei einer Paketgr"o"se von 20 Megabyte ist betr"agt der ausgehende Traffic 8,55 Megabit/s und %
somit auch wieder zwischen den 200 Megabyte und den 2000 Megabyte Paketen. %

\begin{table}
\centering
\begin{tabular}{l%
 r<{\,Mb/s}%
 r<{\,Mb/s}%
}
Tests  					& Eingehender Traffic	& Ausgehender Traffic	\\
\hline
$\diameter$ \hspace{8pt} 20 Megabyte  	& 8,34			& 8,55 			\\
$\diameter$ \hspace{6pt}200 Megabyte  	& 8,42			& 8,67			\\
$\diameter$ 2000 Megabyte	  	& 7,98			& 8,17			\\
\end{tabular}
\caption{Durchschnittlicher eingehender und ausgehender Traffic der einzelnen Testf"alle.}
\label{tab:compTraffic}

\end{table}

Als zweiten vergleichswert wird die Menge an insgesamt verschickten Daten genommen. %
Da es nicht nur das Ziel ist, die Ethernet Ports der Hosts unter der gr"o"stm"oglichen %
Last zu haben sondern das auch effektiv mehr Daten versendet werden sollen, mit denen dann ein anderes %
Hostsystem weiter arbeiten kann ist auch die Menge der verschickten Daten ein entscheidendes Kriterium. Dazu wird die Menge an %
verschickten Daten im Netzwerk miteinander verglichen. Dazu wird in der \cref{tab:compPackages} %
die Anzahl der erfolgreich verschickten Pakete, die verloren gegangen Pakete, die Menge der verschickten %
Byte und der verloren gegangen Byte betrachtet. In der \cref{tab:compPackages} die Ergebnisse aus der %
\cref{tab:VerschickteDaten20Mb}, \cref{tab:VerschickteDaten200Mb} und \cref{tab:VerschickteDaten2000Mb}. %
Bei diesem vergleich ist es offensichtlich welche der Paketgr"o"sen die meisten Pakete verschicken kann, %
hier sieht man das die gr"o"se der Pakete und Anzahl der verschickten Pakete Hand in Hand einhergehen. %
Auch sind vier verlorene Pakete eine sehr geringe Ausfallquote, mit nur 0,04 \% an verloren Paketen, %
ist sind 20 Megabyte Pakete die zweitzuverl"ssigste Paketgr"o"se, was versenden von Daten angeht. %
Am zuverl"assigsten ist in dem Fall der 200 Megabyte Test, in diesem Test ist kein einziges Paket %
verloren gegangen. Anders als im 2000 Megabyte Testfall, bei diesem sind 18 Pakete verloren gegangen, %
was mit 15,56 \% ein im Verh"altnis mit den anderen beiden Tests ein sehr hoher Satz an verlorenen Paketen %
ist. Insgesamt sind bei diesem Test 36 Gigabyte an Daten verloren gegangen. Das ist das 450-fache von dem %
was im 20 Megabyte Test verloren gegangen ist. Damit ist der 2000 Megabyte Testfall nicht der unzuverl"asigste %
von den drei Testf"allen. Auch dieser Testfall schneidet im Vergleich der Menge an verschickten Daten am %
schlechtesten ab, so ist dieser mit 154,30 Gigabyte der Testfall der die wenigsten Informationen %
zwischen den Hosts austauschen konnte. Der Testfall mit den 20 Megabyte Paketen schneidet auch hier %
wieder zwischen dem 2000 Megabyte Test und den 200 Megabyte Test. Dieser hatte insgesamt 213,60 %
Gigabyte "ubertragen. Die meisten "ubertragenen Daten kamen vom 200 Megabyte Testfall mit 224,80 Gigabyte. %

\begin{table}
\centering
\begin{tabular}{l%
 r<{\,}%
   r<{\,}%
r<{\,GB}%
r<{\,}%
}
Tests				& Erfolgreich gesendet			& Erfolglos gesendet 	& Verschickte 	& Verlorene Byte 	\\
\hline
\hspace{8pt} 20 Megabyte  	& 10937					& 4 			& 213,60	& 80 MB		\\
\hspace{6pt}200 Megabyte  	& 1151					& 0			& 224,80	& --		\\
2000 Megabyte	  		& 87					& 18			& 154,30	& 36 GB	\\
\end{tabular}
\caption{Zusammenfassung der Verschickten Pakete und Daten der einzelnen Testf"alle.}
\label{tab:compPackages}
\end{table}

Wie man aus der \cref{tab:compPackages} schlie"sen kann wird im Testfall mit 20 Megabyte Paketen %
die h"ochste Frequenz an verschickten Paketen erreicht mit. Im Testfall mit 200 Megabyte Paketen %
ist Zeitspanne in der ein Paket verschickt wird am zweith"ochsten, gefolgt vom 2000 Megabyte %
Testfall. In der \cref{tab:compZeiten} kann wird genau das aufgezeigt. Dabei wurden die Zeiten %
genommen die in den Logfiles zwischen dem versenden von zwei Paketen gespeichert wurde. Wie man man aus %
dieser Tabelle ablesen kann, wurden durchschnittlich alle 15,80 Sekunden ein 20 Megabyte %
Paket verschickt. Die Standardabweichung f"ur 20 Megabyte Pakete liegt bei 0,11 Sekunden %
woraus man schlie"sen kann das die Pakete zuverl"assig in einem 15,80 $\pm$ 0,11 Sekunden Intervall versendet %
werden. Beim 200 Megabyte Testfall dauerte dies im Schnitt 151,82 Sekunden. Also schon zwei Minuten und 31,82 Sekunden %
was dem 9,61-fachem der Ubertragungsrate des 20 Megabyte Testfalls entspricht. Auch die Standardabweichung ist h"oher %
diese liegt nun bei 3,88 Sekunden, was dem 35,27-fachen der Standardabweichung des 20 Megabyte Testfalls %
entspricht. Betrachtet man nun den 2000 Megabyte Testfall, bemerkt man das die "Ubertragungsdauer eines Paketes %
beim 10,02-fachen des 200 Megabyte Tests liegt oder verglichen mit dem 20 Megabyte Test das 97.10-fache der %
"Ubertragungszeit ben"otigt. Dies spiegelt, sich auch so ungef"ahr in den verschickten Paketen wieder. %
Wie man in der \cref{tab:compPackages} sehen kann entspricht ungef"ahr das 10-Fache der verschickten Pakete im %
2000 Megabyte Test der Anzahl an verschickten Pakete im 200 Megabyte Test. Um das jedoch so zu erreichen m"ussen %
auch die gescheiterten Paket"ubertragungen von dem 2000 Megabyte Test hinzugenommen werden [Wieso dies der Fall ist wird im Abschnitt %
\cref{sec:fehler} erl"autert, Amnk. d. Verf.]. Dasselbe verh"altnis herrscht auch zwischen dem 200 Megabyte Test und %
dem 20 Megabyte Test.  

\begin{table}
\centering
\begin{tabular}{l%
 r<{\,s}%
 r<{\,s}%
 r<{\,s}%
 r<{\,s}%
}
Agent  				& $\diameter $20 Megabyte		& $\diameter $200 Megabyte		& $\diameter $2000 Megabyte		\\
\hline
Dazzle 				& 15,84			 		& 157,65				& 1649,74					\\
Tusk 				& 15,95					& 147,40				& 1445,23				\\
Tinker				& 15,65					& 155,40				& 1574,00				\\
Lion				& 15,75					& 151,82				& 1467,96				\\ 
Agent $\diameter $  		& 15,80					& 153,06			 	& 1534,23				\\   
Agent $\sigma $			& 0,11		 			& 3,88					& 82,52      				\\
\end{tabular}
\caption{Zeiterintervalle bis ein Paket erfolgreich sein Ziel erreicht hat.}
\label{tab:compZeiten}
\end{table}

Als letztes wird der Ping der einzelnen Ger"ate in den verschiedenen Testf"alle betrachtet. %
Obwohl der ICMP Ping auf dem das Programm Pinger aufbaut nicht zuverl"assig ob ein Service %
auf einem Host richtig funktioniert, jedoch gibt es Aufschluss dar"uber ob die Antwort vom eingehende %
ICMP ECHO\_REQUEST stark verz"ogert ist. Sollte die RTT ungew"ohnlich hoch sein, kann es daran liegen %
das der gepingte Host eine sehr hohe CPU-Auslastung hat und den ICMP ECHO\_REQUEST nicht beantworten kann. %
Wie man jedoch aus den \cref{tab:ping20,tab:ping200,tab:ping2000} entnehmen kann ist die Round Trip Time %
sehr gering. Beim 20 Megabyte Testfall Durchschnitt betr"agt die RTT 0,73 ms, w"ahrend beim 200 Megabyte %
Testfall die RTT 0,75 ms betr"agt, womit auch dieser Testfall die hochste RTT hat. Beim 2000 Megabyte Testfall %
liegt die RTT bei 0,73 ms und ist genauso niedrig wie beim 20 Megabyte Testfall. Auff"allig ist jedoch %
das beim Host Lion die h"ochste RTT zu vermessen ist. Betrachtet man die \cref{tab:CPUlastverteilung20Mb,tab:CPUlastverteilung200Mb,tab:CPUlastverteilung2000Mb} %
sieht man das Lion ausser, im 2000 Megabyte Test immer den geringsten Wert im CPU Leerlauf hat. %
Was die hohe RTT von Lion erkl"art. 
  
\begin{table}
\centering
\begin{tabular}{l%
 r<{\,ms}%
 r<{\,ms}%
 r<{\,ms}%
 r<{\,ms}%
 r<{\,ms}%
}
Agent  				& \multicolumn{1}{r}{Dazzle}	& \multicolumn{1}{r}{Tusk}	& \multicolumn{1}{r}{Tinker}	& \multicolumn{1}{r}{Lion}	& \multicolumn{1}{r}{Durchschnitt}		\\
\hline
Dazzle 				& 0,23				& 0,81				& 0,82				& 0,99				& 0,71		\\
Tusk 				& 0,82				& 0,23				& 0,82				& 0,96				& 0,70		\\
Tinker				& 0,83				& 0,81				& 0,23				& 0,99				& 0,71		\\
Lion				& 0,97				& 0,96				& 0,97				& 0,23				& 0,78		\\ 
Durchschnitt			& 0,71				& 0,70				& 0,71 				& 0,79				& 0,73		\\
\end{tabular}
\caption{Die RTT der Pis zueiander beim 20 Megabyte Testfall.}
\label{tab:ping20}
\end{table}

\begin{table}
\centering
\begin{tabular}{l%
 r<{\,ms}%
 r<{\,ms}%
 r<{\,ms}%
 r<{\,ms}%
 r<{\,ms}%
}
Agent  				& \multicolumn{1}{r}{Dazzle}	& \multicolumn{1}{r}{Tusk}	& \multicolumn{1}{r}{Tinker}	& \multicolumn{1}{r}{Lion}	& \multicolumn{1}{r}{Durchschnitt}		\\
\hline
Dazzle 				& 0,23				& 0,85				& 0,84				& 1,00				& 0,73		\\
Tusk 				& 0,85				& 0,24				& 0,85				& 1,02				& 0,74		\\
Tinker				& 0,85				& 0,83				& 0,23				& 1,00				& 0,73		\\
Lion				& 0,99				& 0,98				& 0,99				& 0,23				& 0,80		\\ 
Durchschnitt			& 0,73				& 0,73				& 0,73 				& 0,81				& 0,75		\\
\end{tabular}
\caption{Die RTT der Pis zueiander beim 200 Megabyte Testfall.}
\label{tab:ping200}
\end{table}

\begin{table}
\centering
\begin{tabular}{l%
 r<{\,ms}%
 r<{\,ms}%
 r<{\,ms}%
 r<{\,ms}%
 r<{\,ms}%
}
Agent  				& \multicolumn{1}{r}{Dazzle}	& \multicolumn{1}{r}{Tusk}	& \multicolumn{1}{r}{Tinker}	& \multicolumn{1}{r}{Lion}	& \multicolumn{1}{r}{Durchschnitt}		\\
\hline
Dazzle 				& 0,23				& 0,81				& 0,82				& 0,99				& 0,71		\\
Tusk 				& 0,82				& 0,23				& 0,82				& 0,99				& 0,77		\\
Tinker				& 0,83				& 0,80				& 0,23				& 0,99				& 0,74		\\
Lion				& 0,97				& 0,96				& 0,97				& 0,23				& 0,78		\\ 
Durchschnitt			& 0,71				& 0,70				& 0,71 				& 0,80				& 0,73		\\
\end{tabular}
\caption{Die RTT der Pis zueiander beim 2000 Megabyte Testfall.}
\label{tab:ping2000}
\end{table}

\section{Prozessor Auslastung}
\label{sec:cpulast}
Um die Annahme, die am Ende von \cref{sec:verschickteDaten} gemacht wurde, zu "uberpr"ufen, dass %
ein hohe RTT einer hohen Auslastung der CPU des Zielsystems entpricht, wird nun der Leerlauf den %
die CPUs haben, betrachtet. Da man "uber den Leerlauf R"uckschl"usse "uber die Auslastung der CPU machen kann %
die Theorie die mit der RTT "uber den Host Lion aufgestellt wurde, dass dieser aufgrund der hohen RTT %
unter der h"ochsten Last stand. Wird mit den Werten aus der Tabelle \cref{tab:compCPU} best"agigt. %
Vergleicht man den Leerlauf von Lion mit denen der anderen Hosts, ist der Host Lion derjenige mit %
dem geringsten Leerlauf. Im 2000 Megabyte Testfall jedoch ist die CPU vom Host Tusk am meisten %
ausgelastet. Der Leerlauf beim Host Tusk liegt mit 12,05 \% unter der H"alfte der anderen Hosts im Test. %
Dies ist auch der einzige Test in dem der Host Lion nicht die h"ochste Auslastung hat. Die hohe Auslastung %
des Hosts Tusk h"angt aber auch mit dem Fehler zusammen, der w"ahrend des Tests aufgetreten ist und in dem %
\cref{sec:2000MBTest} und \cref{sec:fehler} schon aufgef"uhrt ist. %

Wie man auch direkt sehen kann sind die drei verbliebenen von den vier Hosts im 2000 Megabyte Test mit einem gro"sen Abstand %
unausgelastet. Mit jeweils 29 \% haben diese Hosts die geringste CPU-Last im Vergleich zu dem 20 Megabyte und %
200 Megabyte Test. Jedoch ist Tusk mit 12,05 \% der ausgelasteste Host in diesem Test. Dies wirkt sich auch auf die %
durchschnittliche Last im Test wieder. Mit 23,67 \% liegt diesmal der 2000 Megabyte Test im Mittelfeld zwischen %
dem 20 Megabyte und 200 Megabyte Test. Bei 20 Megabyte Paketen sind Hosts im Test mit 23,83 \% die mit dem %
h"ochsten Leerlaufprozess. W"ahrend beim 200 Megabyte Test mit \mbox{23,35 \%} der geringste Leerlaufprozess erreicht %
ist. Aus der Standardabweichung der einzelnen Tests kann man auch schlussfolgern, welcher der Tests eine m"oglichst %
gleichm"a"sige Belastung der Hosts hat. So werden beim 20 Megabyte Test, alle Hosts mit einer Standardabweichung %
von $\pm$ 0,10 \% am gleichm"a"sigsten belastet. Gefolgt vom 200 Megabyte Test. Bei diesem Test ist auch noch eine %
geringe Standardabweichung gegeben mit $\pm$ 0,98 \% sind. Auch in diesem Test sind die Hosts gleichm"a"sig belastet. %
Anders als beim 2000 Megabyte Test ist die CPU Lastenverteilung ist sehr ungleich. Mit $\pm$ 7,05 \% ist die Lastverteilung %
am ungleichm"a"sigsten. Der Host Tusk hat in diesem Test die h"ochste Last von allen Host. % 


\begin{table}
\centering
\begin{tabular}{l%
 r<{\,\%}%
 r<{\,\%}%
 r<{\,\%}%
}
Agent  				& \multicolumn{1}{r}{20 Megabyte Idle}			& \multicolumn{1}{r}{200 Megabyte Idle}			& \multicolumn{1}{r}{2000 Megabyte Idle}		\\
\hline
Dazzle 				& 25,45			 				& 21,83							& 29,48					\\
Tusk 				& 23,26							& 25,57							& 12,05					\\
Tinker				& 23,78							& 24,55							& 29,16					\\
Lion				& 22,82							& 21,48							& 29,16					\\ 
Agent $\diameter $  		& 23,83							& 23,35					 		& 23,67					\\   
Agent $\sigma $			& 0,10		 					& 0,98							& 7,05      				\\
\end{tabular}
\caption{Leerlauf der CPUs im vergleich. Werte aus den \cref{tab:CPUlastverteilung20Mb,tab:CPUlastverteilung200Mb,tab:CPUlastverteilung2000Mb}.}
\label{tab:compCPU}
\end{table}


\section{Festplatten Auslastung}
\label{sec:festlast}

In diesem Abschnitt wird die Auslastung der Festplatten betrachtet. %
Dabei wird die Metrik der Input Output Operations per Second hinzugezogen. %
Diese gibt an wieviele Lese oder Schreib Operationen in der Sekunde auf der Festplatte %
stattfinden. Je h"oher der Wert ist desto mehr wird auf der Festplatte gelesen oder geschrieben. %
Daraus folgt, jedoch auch das die CPU eine h"ohere I/O Wait Time hat und somit l"anger darauf %
wartet das die Festplatte mit einer Schreib oder Lese Operation abgeschlossen hat. %
Deshalb wird als erstes die CPU I/O Wait Time betrachtet. Wie man aus der \cref{tab:compIoWaitCpu} %
sehen kann ist beim 20 Megabyte Test die geringste Wartezeit mit 1,26 \% f"ur die CPU. Danach gefolgt von dem 200 Megabyte %
Testfall, dieser hat mit 5,03 \% die zweith"ochste I/O Wait Time. Beim 2000 Megabyte Test ist mit 5,85 \% die %
CPU am l"angsten damit besch"aftigt auf den Abschluss einer Lese oder Schreib Operation zu warten. Jedoch ist %
die Differenz zwischen dem 200 und 2000 Megabyte Test nicht so signifikant wie zwischen dem 20 Megabyte Test. %

\begin{table}
\centering
\begin{tabular}{l%
 r<{\,\%}%
 r<{\,\%}%
 r<{\,\%}%
}
Agent	  			& \multicolumn{1}{r}{20 Megabyte}	 	& \multicolumn{1}{r}{200 Megabyte}		& \multicolumn{1}{r}{2000 Megabyte} 		\\	
\hline
Dazzle 				& 1,42						& 7,02						& 7,33	        \\
Tusk 				& 0,22						& 1,20						& 1,27		\\
Tinker				& 1,65						& 5,67						& 8,11	 	\\
Lion				& 1,73						& 6,68						& 6,68	 	\\
Agent $\diameter $  		& 1,26						& 5,14						& 5,85		\\   
Agent $\sigma $ 		& 0,20 						& 2,33						& 2,69			\\
\end{tabular}
\caption{CPU Wartezeit auf den Abschluss einer Lese oder Schreib Operation, Werte aus den \cref{tab:CPUlastverteilung20Mb,tab:CPUlastverteilung200Mb,tab:CPUlastverteilung2000Mb}.}
\label{tab:compIoWaitCpu}
\end{table}

In der \cref{tab:compOps} sind die jeweiligen Lese und Schreib Operationen pro Sekunde der einzelen %
Testf"alle aufgelistet. Wie man sieht geschehen die geringsten I/O Operationen w"ahren dem 20 Megabyte %
Test. Das spiegelt sich auch in der I/O Wait Time der CPU wieder, mit durchschnittlich 120,73 Ops/s %
wird bei diesem Test am wenigstens geschrieben. Wie die \cref{tab:NormalbetriebIoStat20Mb} auch zeigt %
wird in diesem Testfall gar keine Daten gelesen. Das erkl"art auch wieso die Werte so f"ur die Lese %
und Schreib Operationen so gering sind. W"ahrend erst beim 200 Megabyte Test die ersten Lese Operationen %
stattfinden mit durchschnittlich 907,56 Byte/s wie man aus \cref{tab:NormalbetriebIoStat200Mb} lesen kann. Dies erkl"art auch den Anstieg der Lese und Schreib Operationen %
Die der Durchschnitt der I/O Operationen/s zwischen den Tests mit 200 Megabyte Paketen und 2000 Megabyte Paketen %
liegt nahe beinander. Dasselbe ist auch beim Durchschnitt der CPU I/O Wait Time zu beobachten, jedoch %
haben sich die Verh"altnisse ver"andert. W"ahrend beim 200 Megabyte Test mehr geschrieben als gelesen wurde %
ist dies beim Testfall mit 2000 Megabyte andersrum, in 2000 Megabyte Testfall wird mit durchschnittlich 2,30 KB/s
gelesen nach \cref{tab:NormalbetriebIoStat2000Mb}. Daf"ur ist die durchschnittliche Menge die geschrieben wird %
auf 2,12 KB/s, w"ahrend die vom 200 Megabyte Test noch bei 2,65 KB/s lag.   


\begin{table}
\centering
\begin{tabular}{l%
 r<{\,Ops/s}%
 r<{\,Ops/s}%
 r<{\,Ops/s}%
}
Agent	  			& \multicolumn{1}{r}{20 Megabyte}	 	& \multicolumn{1}{r}{200 Megabyte}		& \multicolumn{1}{r}{2000 Megabyte} 		\\	
\hline
Dazzle 				& 153,09					& 263,34					& 257,88	        \\
Tusk 				& 42,40						& 122,64					& 132,75		\\
Tinker				& 142,50					& 234,48					& 257,25	 	\\
Lion				& 144,31					& 248,08					& 248,08	 	\\
Agent $\diameter $  		& 120,73					& 218,39					& 223,99		\\   
Agent $\sigma $ 		& 45,31 					& 56,58						& 52,82			\\
\end{tabular}
\caption{Lese und Schreibzugriffe auf die Festplatten w"ahrend der Testl"aufe, Werte aus den \cref{tab:NormalbetriebIoStat20Mb,tab:NormalbetriebIoStat200Mb,tab:NormalbetriebIoStat2000Mb}.}
\label{tab:compOps}
\end{table}

\section{Aufgetretene Fehler}
\label{sec:fehler}

Wie schon in \cref{sec:2000MBTest} erw"ahnt kam es bei diesem Test zu einem Fehler %
bei einem der Hostsysteme. Alle Pakete, die zum Host verschickt werden sollten, %
haben ihr Ziel nicht erreicht. Dadurch sind sehr viele Pakete verloren gegangen. %
Insgesamt konnte vom Host Tusk kein einziges Paket angenommen werden. So sind alle %
fehlerhaften Pakete auf ihn zur"uckzuf"uhren. Das Monitoring Tool Zabbix hat dazu selber %
eine Warnung ausgegeben. Diese besagt, dass der freie Festplattenspeicher zu mehr als 80 \% %
belegt ist. Dies gibt einen ersten Hinweis darauf, dass die 2000 Megabyte Dateien an %
diesen Host nicht richtig gesendet werden k"onnen.  

\begin{figure}[htbp]
\centering
\includegraphics*[width=1\linewidth]{Abb/Tusk20prozent}

\caption{Fehlermeldung bez"uglich der Festplatte von Tusk}
\label{fig:Eth0DazzleStandard}
\end{figure}

Schaut man in die Logfiles und betrachtet die Zeit die zwischen dem Versenden von zwei Paketen vergeht sieht %
man, dass die Zeit die bis zum versenden des n"achsten Packetes %
sehr hoch ist in diesem Fall 17 Minuten. 

\begin{verbatim}
[Thu 16 Jun 19:36:11 CEST 2016]TRANSPORT FAILED 
[Thu 16 Jun 19:53:53 CEST 2016]TRANSPORT SUCCESSFUL
\end{verbatim} 

Das "Ubertragen von erfolgreichen Paketen dauert nach der \cref{tab:compZeiten} durchschnittlich 1534,23 Sekunden %
also ungef"ahr 25 Minuten. Mit diesem Wissen kann man die Schlussfolgerung ziehen, dass der Abbruch des verschickten %
Paketes an den Host Tusk erst sehr sp"at in der "Ubertragung passiert. SCP "uberpr"uft vor dem Versenden einer Datei nicht %
nach, ob genug Speicherplatz f"ur diese Datei vorhanden ist. So kommt es auch, dass der Host Tusk auch im %
Test auf dem Ethernet Port eine eingehende Last aufweisen kann. Da SCP auf der SSH aufbaut welches das %
TCP Protokoll verwendet, werden die verschickten Dateien segmentiert. F"ur gew"ohnlich sind solche Segmente %
1500 Byte gro"s. Dies bedeutet, dass der Host Tusk solange 1500 Byte TCP-Segmente empf"angt, bis bei diesem %
die Festplatte voll ist und keine weiteren Segmente mehr annehmen kann. Daraus folgt dann, dass die Daten"ubertragung %
abgebrochen wird, Zabbix hat f"ur die verschickten Segmente Statistiken sammeln k"onnen, wie z.B: Den eingehenden %
Traffic und Belastung der CPU, obwohl der Host nie die komplette Datei empfangen hat. % 

\section{Schlussfolgerung}
\label{sec:schluss}

Wie bei vielem in der Informatik ist es nicht m"oglich \emph{das Beste} zu definieren %
und danach eine Auswahl zu treffen. Es muss st"andig abgewogen werden was in einem %
Use Case am wichtigsten ist. Jeder der gezeigten Tests, hat den anderen %
Tests gegen"uber Vor- und Nachteile. Deshalb betone ich an dieser Stelle nochmal die %
Zielsetzung die in \cref{cha:einfuehrung} definiert wurde. 

\begin{quotation}
Das Ziel dieser Bachelorarbeit ist es herauszufinden, ob "uber die Paketgr"o"se eine %
Optimierung der Performanz der Endger"ate und der Daten"ubertragung im Netzwerk m"oglich ist. %
Dabei wird versucht die gr"o"stm"ogliche Menge an Daten im Netzwerk auszutauschen ohne die %
Funktionalit"at eines der Hosts zu gef"ahrden. %
\end{quotation}

Wie man in \cref{cha:versuche} sehen konnte, gibt es signifikante Unterschiede bei den in dieser %
Ausarbeitung durchgef"uhrten Testf"alle. Deshalb wird der Fokus erstmal darauf liegen, die %
Anforderung der Zielsetzung zu erf"ullen. Mit dem Abschnitt der Zielsetzungi, das die %
Funktionalit"at der Hosts nicht gef"ahrdet wird, f"allt der 2000 Megabyte Testversuch %
schonmal als eine geeignete Paketgr"o"se raus. Aufgrund der in \cref{sec:fehler} % 
hohen Paketgr"o"se kann der Host Tusk nicht fehlerfrei arbeiten. Seine CPU %
ist zu 88 \% ausgelastet und Tusk kann keine empfangen. Dies alleine ist schon %
ein Kriterium, dass 2000 Megabyte Pakete ausschlie"st. Dass auch noch 35,16 Gigabyte %
an Daten im Netz verloren gegangen und insgesamt nur 154,30 Gigabyte verschickt wurden, %
erf"ullt auch die zweite Anforderung nicht, n"amlich m"oglichst viele Daten zu verschicken. %
Damit f"allt die m"ogliche Auswahl auf die 20 Megabyte Dateien oder die 200 Megabyte %
Dateien. Dazu wird als erstes die Menge der verschickten Gigabyte miteianander verglichen. %
Aus der \cref{tab:compPackages} l"asst sich ablesen, dass bei 200 Megabyte 224,80 %
Gigabyte an Daten verschickt worden sind und keine einziges Packet verloren gegangen ist. %
Beim 20 Megabyte Test wurden hingegen nur 213,60 Gigabyte verschickt. Dabei sind 80 Megabyte %
an Daten verloren gegangeni, was viel weniger als 0,01 \% ist. Deshalb sind 200 Megabyte Testdateien %
f"urs erste am geeigneitsten. Jedoch muss noch der zweite Punkt beachtet werden, n"amlich % 
dass die Hosts im System nicht komplett "uberlastet werden und noch simple Aufgaben %
erf"ullen k"onnen. Dazu werden die Werte aus der \cref{tab:compCPU} miteinanander verglichen. %
Der Leerlaufprozess beim 20 Megabyte Test betr"agt 23,83 \% und beim 200 Megabyte Test \% 23,35 \%. %
Obwohl das hei"st das, dass die Hosts im Test mit 20 Megabyte Dateien weniger belastet waren als im %
Test mit 200 Megabyte Dateien, ist der Unterschied sehr geringf"ugig. Deshalb h"alt man sich %
bei der Entscheidung, welche der Dateigr"o"sen am besten geeignet ist daran, welche der beiden %
Testf"alle die meisten Daten verschickt hat, welche der 200 Megabyte Test war. %

Da aber, wie am Anfang dieses Abschnittes erw"ahnt, die optimale Dateigr"o"se %
vom jeweiligen Use Case abh"angig ist muss man von Use Case zu Use Case unterscheiden. %
Als Beispiel in dem die 20 Megabyte Datei von Vorteil w"are, ist in eine Netzwerkanwendung %
wo die Zeit ein wichtige Rolle spielt und es nicht auf Menge der Daten ankommt sondern %
das schnell Daten ankommen die 20 Megabyte Datei gr"o"se die bessere Wahl. %




%----------------- KAPITEL : FAZIT  ----------------- %	
%\chapter{Fazit}
\chapter{Anhang}
\label{cha:anhang}

\section{Test Ruhezustand}

\begin{figure}[htbp]
\centering
\includegraphics*[width=1\linewidth]{Abb/ZabbixDazzle/NoTraffic/DazzleNoTrafficEth0}

\caption{Traffic auf Eth0 ohne Hintergrundrauschen auf Dazzle.}
\label{fig:Eth0DazzleNotraffic}
\end{figure}


\begin{figure}[htbp]
\centering
\includegraphics*[width=0.9\linewidth]{Abb/ZabbixDazzle/NoTraffic/DazzleNoTrafficIoStat}

\caption{I/O Statistik von Dazzle ohne Hintergrundrauschen.}
\label{fig:IoStatDazzleNotraffic}
\end{figure} %

\begin{figure}[htbp]
\centering
\includegraphics*[width=0.9\linewidth]{Abb/ZabbixDazzle/NoTraffic/DazzleNoTrafficCPU}

\caption{Prozentuale Lastverteilung auf der CPU von Dazzle, ohne Hintergrundrauschen.}
\label{fig:CPUDazzleNotraffic}
\end{figure}
\FloatBarrier
\section{Test 20 Megabyte Tests}

\begin{figure}[htbp]
\centering
\includegraphics*[width=0.9\linewidth]{Abb/ZabbixDazzle/Standard/Standard20}
\caption{Traffic auf Eth0 bei 20 Megabyte Dateien auf Dazzle.}
\label{fig:Standard20}
\end{figure}


\begin{figure}[htbp]
\centering
\includegraphics*[width=0.9\linewidth]{Abb/ZabbixDazzle/Standard/IoStatStandard20}

\caption{I/O Statistik von Dazzle beim 20 Megabyte Betrieb.}
\label{fig:IoStatStandard20}
\end{figure} %

\begin{figure}[htbp]
\centering
\includegraphics*[width=0.9\linewidth]{Abb/ZabbixDazzle/Standard/CPUStandard20}

\caption{Prozentuale Lastverteilung auf der CPU bei 20 Megabyte Paketen von Dazzle.}
\label{fig:CPUDazzleStandard}
\end{figure}
\FloatBarrier
\newpage
\section{200 Megabyte Tests}

\begin{figure}[htbp]
\centering
\includegraphics*[width=0.9\linewidth]{Abb/ZabbixDazzle/Standard/Standard200}
\caption{Traffic auf Eth0 bei 200 Megabyte Dateien auf Dazzle.}
\label{fig:Standard200}
\end{figure} %


\begin{figure}[htbp]
\centering
\includegraphics*[width=0.9\linewidth]{Abb/ZabbixDazzle/Standard/IoStatStandard200}

\caption{I/O Statistik von Dazzle beim 200 Megabyte Betrieb.}
\label{fig:IoStatDazzleStandard200}
\end{figure} %

\begin{figure}[htbp]
\centering
\includegraphics*[width=0.9\linewidth]{Abb/ZabbixDazzle/Standard/CPUStandard200}

\caption{Prozentuale Lastverteilung auf der CPU bei 200 Megabyte Paketen von Dazzle.}
\label{fig:CPUDazzleStandard200}
\end{figure}
\FloatBarrier
\newpage
\section{2000 Megabyte Pakete}

\begin{figure}[htbp]
\centering
\includegraphics*[width=0.9\linewidth]{Abb/ZabbixTusk/Standard/Standard2000}
\caption{Traffic auf Eth0 bei 2000 Megabyte Dateien auf Dazzle.}
\label{fig:Standard2000}
\end{figure} %


\begin{figure}[htbp]
\centering
\includegraphics*[width=0.9\linewidth]{Abb/ZabbixDazzle/Standard/IoStatStandard2000}

\caption{I/O Statistik von Dazzle beim 2000 Megabyte Betrieb.}
\label{fig:IoStatDazzleStandard2000}
\end{figure} %

\begin{figure}[htbp]
\centering
\includegraphics*[width=0.9\linewidth]{Abb/ZabbixDazzle/Standard/CPUStandard2000}

\caption{Prozentuale Lastverteilung auf der CPU bei 2000 Megabyte Paketen von Dazzle.}
\label{fig:CPUDazzleStandard2000}
\end{figure}

\section{Inhalt der CD}
Auf der CD sind folgende Inhalte:
\begin{itemize}
\item Die Logfiles zu den Tests, 
\item die selbsterstellten Programme,
\item die Graphen der restlichen Hosts,
\item eine Libre Office Calc Tabelle.
\end{itemize}




















%----------------- KAPITEL : LITERATUR  ----------------- %	
%\printbibliography
%----------------- KAPITEL : ANHANG  ----------------- %	


\end{document}
